% The encodings sec:vartag compares, drawn as the 64-bit word a slot holds.
% McKenney asked for a diagram of this section's nomenclature: tag width and
% alignment, low versus high bits, the reader's test, recovering the proxy by
% subtraction, and the zero-anchored reserved range with its index handles are
% all defined in prose only.
%
% Drawn in TikZ (not an imported image) so the figure ships as source in the
% arXiv submission and needs no external toolchain to rebuild.  Colors follow
% fig:anatomy: orange for the slot's facility-owned bits, blue for a proxy.
%
% The picture must make one thing obvious per row: where the marker sits in the
% word, what the reader's test is, and how the proxy is recovered.
\begin{figure}[tb]
\centering
\begin{tikzpicture}[
  font=\small,
  fld/.style  = {draw, minimum height=6.5mm, inner sep=0pt, font=\footnotesize},
  addr/.style = {fld, fill=gray!6},
  paddr/.style= {fld, fill=blue!6},
  tagf/.style = {fld, fill=orange!22, very thick, font=\footnotesize\ttfamily},
  lowf/.style = {fld, fill=gray!15, font=\footnotesize},
  bit/.style  = {font=\scriptsize\ttfamily, text=black!60},
  ttl/.style  = {font=\footnotesize\bfseries, anchor=west},
  lbl/.style  = {font=\footnotesize\itshape, anchor=east, text=black!75},
  op/.style   = {font=\footnotesize, anchor=west, align=left},
]
% Geometry: a word is 7.0 cm wide; a 4-bit field is 1.3 cm.
\def\W{7.0}\def\L{1.3}\def\X{2.2}  % word width, 4-bit field width, word x-origin
\def\O{9.45}                        % x of the annotation column

% ---------------------------------------------------------------- (a) low bits
\node[ttl] at (0, 0.55) {(a) Low-bit tag --- this facility};
% plain value
\node[lbl] at (\X-0.1, 0) {plain value};
\node[addr, minimum width={(\W-\L)*1cm}, anchor=west] (a1) at (\X, 0)
  {a node address, or any value};
\node[lowf, minimum width={\L cm}, anchor=west] at (a1.east) {$\neq$ tag};
% parked
\node[lbl] at (\X-0.1, -0.8) {parked slot};
\node[paddr, minimum width={(\W-\L)*1cm}, anchor=west] (a2) at (\X, -0.8)
  {proxy address, 16-byte aligned};
\node[tagf, minimum width={\L cm}, anchor=west] (a2t) at (a2.east) {tag};
% bit numbers
\node[bit, anchor=north west] at ($(a2.south west)+(0,-0.02)$) {63};
\node[bit, anchor=north east] at ($(a2.south east)+(0,-0.02)$) {4};
\node[bit, anchor=north west] at ($(a2t.south west)+(0,-0.02)$) {3};
\node[bit, anchor=north east] at ($(a2t.south east)+(0,-0.02)$) {0};
\node[op] at (\O, -0.4)
  {test: \code{(v \& tag) == tag}, one AND\\
   proxy: \code{v - tag}, folded into the load};

% --------------------------------------------------------------- (b) high bits
\node[ttl] at (0, -2.05) {(b) High-bit tag --- the alternative};
\node[lbl] at (\X-0.1, -2.6) {parked slot};
\node[tagf, minimum width={\L cm}, anchor=west] (b1t) at (\X, -2.6) {tag};
\node[paddr, minimum width={(\W-\L)*1cm}, anchor=west] (b1) at (b1t.east)
  {proxy address};
\node[bit, anchor=north west] at ($(b1t.south west)+(0,-0.02)$) {63};
\node[bit, anchor=north east] at ($(b1t.south east)+(0,-0.02)$) {60};
\node[bit, anchor=north east] at ($(b1.south east)+(0,-0.02)$) {0};
\node[op] at (\O, -2.6)
  {test: the mask needs a register\\
   proxy: an extra instruction first};

% ------------------------------------------------------------ (c) reserved range
\node[ttl] at (0, -3.85) {(c) Reserved range, anchored at zero --- no bits spent};
\node[lbl] at (\X-0.1, -4.4) {parked slot};
\node[addr, minimum width={(\W-2.6)*1cm}, anchor=west] (c1) at (\X, -4.4)
  {zero: below any allocation};
\node[tagf, minimum width={2.6cm}, anchor=west, font=\footnotesize] (c1t) at (c1.east)
  {index $<N$};
\node[bit, anchor=north west] at ($(c1.south west)+(0,-0.02)$) {63};
\node[bit, anchor=north east] at ($(c1t.south east)+(0,-0.02)$) {0};
\node[op] at (\O, -4.4)
  {test: \code{v < N}, one compare\\
   proxy: \code{pool[v]}};
\end{tikzpicture}
\caption{The encodings \cref{sec:vartag} compares, on a 64-bit slot. (a) While
a commit is in flight this facility parks a proxy's address in the slot with the
tag added in its low bits, which the proxy's 16-byte alignment guarantees clear;
every value the slot may otherwise hold must not carry the tag pattern
(\cref{sec:tagcontract}). Subtracting the tag recovers the same address as
masking it off, since the proxy's own tag bits are clear. (b) A high-bit tag
leaves the low bits to the embedder; on x86-64 a mask of several high bits does
not fit an immediate, and recovering the proxy takes an instruction of its own.
(c) A reserved range spends no bits of the word: a handle is an index below the
pool size $N$, which must cover every proxy live within a grace period.}
\label{fig:tagword}
\end{figure}
